%% literature_review_chapters/ch2_5_literature_review.tex
%%
%% Standalone literature-review draft covering chapters 2--5 of the thesis.
%% Not yet wired into thesis/main.tex. Notation reuses thesis/macros.tex.
%%
%% Structure:
%%   Chapter 2: Continuous Foundations of Diffusion Models
%%   Chapter 3: Discrete Diffusion Models
%%   Chapter 4: Masked Diffusion Language Models
%%   Chapter 5: Remasking, Self-Correction, and Informed Correctors
%%
%% Audience: mathematically mature ML reader. The chapters foreshadow but do
%% not develop the contribution-chapter framework (G(S), A(S), Q(S),
%% Theorem A/B/B', Diagnostic Framework C); those remain in ch6.

%% =====================================================================
%% Chapter 2 -- Continuous Foundations of Diffusion Models
%% =====================================================================

\chapter{Continuous Foundations of Diffusion Models}
\label{lit:ch:continuous-diffusion}

This chapter reviews continuous-state diffusion models at the level needed to
motivate the discrete masked setting that follows. The aim is not a complete
derivation of the score-based framework. It is to isolate the structural
ingredients that survive the move to text: a designer-chosen forward
corruption process, a learned reverse process, a notion of score or denoising
target, an inference-time discretisation, and the practice of inserting
auxiliary correction transitions along an otherwise fixed sampling
trajectory. The last of these is the conceptual root of corrector scheduling
and is introduced here in its native continuous form before it is recast for
discrete states in
Chapters~\ref{lit:ch:discrete-diffusion}--\ref{lit:ch:correctors}.

\section{From likelihood-based generative modelling to diffusion}
\label{lit:sec:gen-modeling}

Generative modelling on a data space $\mathcal{X}$ seeks a parametric
distribution $p_\theta$ that minimises a divergence to the unknown data
distribution $\pdata$. Maximum-likelihood estimation,
$\theta^\star \in \argmin_\theta \KL{\pdata}{p_\theta}$, is attractive because
it is consistent and gradient-friendly when $\log p_\theta$ is tractable, but
tractable density families are restrictive. Autoregressive models exploit the
chain rule but are intrinsically sequential at sampling. Normalising flows
admit exact likelihood through bijective parameterisations at the cost of
architectural restrictions. Variational autoencoders relax the likelihood
through a stochastic latent and an evidence lower bound but introduce an
amortisation gap.

Diffusion models occupy a different position in this design space. They
factor generation through a long Markov chain of conditionally tractable
transitions, with the high-noise endpoint chosen to be a reference
distribution from which exact sampling is trivial. The construction is due
to \textcite{sohldickstein2015deep}, who proposed parameterising both a
forward Markov noising chain and a learned reverse Markov chain and training
the latter to maximise a variational lower bound on the marginal data
log-likelihood. The construction was largely dormant until
\textcite{ho2020ddpm} simplified the training loss into a noise-prediction
regression on Gaussian forward chains. In parallel, \textcite{song2019ncsn}
introduced score-based generative models via denoising score matching across
multiple noise scales. \textcite{song2021score} unified these viewpoints
through a continuous-time stochastic differential equation. The literature
has since converged on a small set of equivalent formulations that differ in
parameterisation and training objective but agree on the geometry of the
underlying process.

The relevance for the present thesis is structural rather than
algorithmic. Diffusion models exhibit a consistent split between an inference
trajectory the designer fixes (a sequence of intermediate noisy states) and
a learned per-step transition the network supplies. Inference-time
interventions---accelerated solvers, corrector steps, guidance---operate on
the trajectory, not on the trained network. This split survives the move
from $\mathbb{R}^d$ to $\mathcal{V}^L$ even when nearly all of the
mathematical machinery of the continuous theory does not.

\section{Forward noising and reverse denoising}
\label{lit:sec:fwd-rev}

Fix a forward noising process indexed by $t\in[0,1]$ (or, in discrete
time, $t\in\{1,\ldots,T\}$) with kernel $q(x_t\mid x_0)$ and marginals
$q_0 = \pdata$ and $q_T \approx p_{\mathrm{ref}}$, where
$p_{\mathrm{ref}}$ is a reference distribution from which exact sampling
is trivial---typically $\mathcal{N}(0,I)$ or
$\mathcal{N}(0,\sigma_{\max}^2 I)$. The reverse process
$p_\theta(x_{t-1}\mid x_t)$ is parameterised so that, when initialised at
$x_T \sim p_{\mathrm{ref}}$ and composed iteratively, its marginal at
$t=0$ approximates $\pdata$.

The DDPM-style discrete-time forward chain uses a variance schedule
$\{\betanoise_t\}_{t=1}^{T}\subset(0,1)$ and the kernel
\[
  q(x_t \mid x_{t-1})
  = \mathcal{N}\!\big(x_t;\, \sqrt{1-\betanoise_t}\,x_{t-1},\, \betanoise_t I\big).
\]
A direct calculation gives the closed-form $t$-step kernel,
\[
  q(x_t \mid x_0)
  = \mathcal{N}\!\big(x_t;\, \sqrt{\abardpm_t}\,x_0,\, (1-\abardpm_t)I\big),
  \qquad
  \abardpm_t = \prod_{s=1}^{t}(1-\betanoise_s).
\]
The shortcut from $x_0$ to $x_t$ is what enables stochastic-gradient
training: one samples a clean datum $x_0\sim\pdata$, a time index $t$, and
noise $\varepsilon\sim\mathcal{N}(0,I)$, and forms
$x_t = \sqrt{\abardpm_t}\,x_0 + \sqrt{1-\abardpm_t}\,\varepsilon$ in one
step.

The exact reverse conditional $q(x_{t-1}\mid x_t,x_0)$ is also Gaussian with
closed-form mean and variance, but it depends on the unobserved $x_0$. The
model therefore predicts a quantity from which the reverse mean can be
reconstructed: the noise $\varepsilon$, the clean datum $x_0$, or the
marginal score $\nabla_{x}\log q_t(x_t)$. The three choices are equivalent
up to a time-dependent reweighting and produce the same minimiser
\parencite{ho2020ddpm,song2021score}.

\section{Denoising diffusion probabilistic models}
\label{lit:sec:ddpm}

\textcite{ho2020ddpm} train an $\varepsilon$-prediction network
$\etheta(x_t,t)$ to minimise the simple regression loss
\[
  \mathcal{L}_{\mathrm{simple}}
  = \E_{t,x_0,\varepsilon}\,
    \big\|\varepsilon
        - \etheta\!\big(\sqrt{\abardpm_t}\,x_0
                       + \sqrt{1-\abardpm_t}\,\varepsilon,\, t\big)\big\|^2,
\]
where $t$ is sampled uniformly on $\{1,\ldots,T\}$ and
$\varepsilon\sim\mathcal{N}(0,I)$. This loss is a reweighted version of
the variational lower bound on $\log p_\theta(x_0)$, and \textcite{ho2020ddpm}
report that the unweighted form performs better empirically. The
corresponding ancestral sampler is
\[
  x_{t-1}
  = \tfrac{1}{\sqrt{1-\betanoise_t}}\Big(x_t
       - \tfrac{\betanoise_t}{\sqrt{1-\abardpm_t}}\,\etheta(x_t,t)\Big)
    + \sigma_t z,
  \qquad z\sim\mathcal{N}(0,I),
\]
with a schedule-dependent variance $\sigma_t$. The combined effect is that
the network supplies the predicted denoising direction and the designer
supplies a Gaussian fluctuation calibrated to the local noise level.

The objective admits several equivalent reformulations. With
$\xhat_\theta(x_t,t)
 = (x_t - \sqrt{1-\abardpm_t}\,\etheta(x_t,t))/\sqrt{\abardpm_t}$ the same
network is an $x_0$ predictor; with $\stheta(x_t,t) =
-\etheta(x_t,t)/\sqrt{1-\abardpm_t}$ it is a score estimator. These
equivalences are routinely used in practice to switch between
parameterisations without retraining.

The DDPM training loss separates the network from the schedule. The
network learns a single function over noisy inputs and time indices; the
schedule $\{\betanoise_t\}$ is fixed in advance and determines the marginals
along which both training samples and reverse trajectories pass. As long as
the chain is long enough to make $q(x_T)\approx\pi$, the schedule details
affect the difficulty of the regression at each $t$ but not the existence of
a Bayes-optimal solution. This decoupling is preserved in masked
diffusion, where the schedule fixes a marginal mask ratio and the network
learns to denoise across the implied family of partially masked sequences.

\section{Score matching and the score-based view}
\label{lit:sec:score-matching}

The score function of a smooth distribution $p$ on $\mathbb{R}^d$ is
$s(x) = \nabla_x\log p(x)$. \textcite{hyvarinen2005estimation} introduced
score matching as a way to fit unnormalised densities without computing a
partition function: under regularity conditions, the implicit objective
\[
  \mathcal{J}_{\mathrm{ISM}}(\theta)
  = \E_{x\sim p}\,\Big[
      \tfrac{1}{2}\|\stheta(x)\|^2
      + \tr\!\big(\nabla_x \stheta(x)\big)
    \Big]
\]
agrees, up to a constant, with the squared-error score-matching loss. The
trace term is computationally inconvenient at scale.

\textcite{vincent2011connection} showed that score matching against a
smoothed distribution $p_\sigma$, obtained by convolving $p$ with
$\mathcal{N}(0,\sigma^2 I)$, is equivalent to denoising score matching:
\[
  \mathcal{J}_{\mathrm{DSM}}(\theta;\sigma)
  = \E_{x_0\sim p,\,\tilde x\sim\mathcal{N}(x_0,\sigma^2 I)}\,
    \big\|\stheta(\tilde x) - \nabla_{\tilde x}\log p(\tilde x\mid x_0)\big\|^2.
\]
Because $\nabla_{\tilde x}\log p(\tilde x\mid x_0) = -(\tilde x-x_0)/\sigma^2$
when the conditional is Gaussian, the inner gradient has a closed form and
the loss is a simple regression. \textcite{song2019ncsn} train a
noise-conditional score network across a geometric sequence of noise scales
$\{\sigma_i\}_{i=1}^{N}$ and sample by annealed Langevin dynamics. Subsequent
work refines training stability and sampling \parencite{song2020improvedncsn}.

The score-based view does not commit to a normalisable density. It instead
fits a vector field $\stheta(x,t)$ that approximates $\nabla_x\log q_t(x)$
across time. Sampling becomes the integration of a stochastic process whose
drift is built from this vector field. This is structurally different from
autoregressive or flow-based models: there is no exact likelihood evaluation,
but sampling cost is decoupled from the likelihood computation and can be
tuned independently.

\section{Continuous-time diffusion via stochastic differential equations}
\label{lit:sec:sde}

\textcite{song2021score} unify DDPM and NCSN by writing the forward process
as a stochastic differential equation
\[
  \dif x_t = f(x_t,t)\,\dif t + g(t)\,\dif \Wt,
  \qquad x_0 \sim \pdata,
\]
where $\Wt$ is a standard Brownian motion. Variance-preserving (VP) and
variance-exploding (VE) SDEs correspond to the DDPM and NCSN constructions
respectively. The marginal density $q_t$ is governed by a Fokker--Planck
equation. The classical reverse-time formula \parencite{anderson1982reverse}
gives the reverse SDE
\[
  \dif x_t = \big[f(x_t,t) - g(t)^2\,\nabla_x\log q_t(x_t)\big]\,\dif t
              + g(t)\,\dif \bar \Wt,
\]
where $\bar\Wt$ is a Brownian motion under reverse time. Substituting a
learned score $\stheta(x_t,t)\approx\nabla_x\log q_t(x_t)$ yields a sampling
SDE. \textcite{song2021score} additionally introduce the probability-flow ODE
\[
  \dif x_t = \big[f(x_t,t) - \tfrac{1}{2}g(t)^2\,\stheta(x_t,t)\big]\,\dif t,
\]
which has the same time marginals as the reverse SDE but is deterministic
conditional on $x_T$.

The continuous-time view is conceptually clean. Training estimates a single
score field $\stheta(\cdot,\cdot)$; sampling integrates either the reverse
SDE or the ODE; the choice of forward SDE family becomes a design variable
rather than an incidental discretisation. The view also exposes sampling
cost as a function of integrator and step count. Subsequent literature
exploits this exposed structure (\S\ref{lit:sec:fast-samplers}).

The SDE formulation is also where the connection to continuous-time Markov
chains on discrete state spaces becomes clearest. A continuous-state SDE
generates a Markov process whose generator acts on smooth test functions; a
continuous-time Markov chain on a finite alphabet generates a Markov process
whose generator is a rate matrix. Discrete diffusion models, reviewed in
Chapter~\ref{lit:ch:discrete-diffusion}, replace the diffusion operator by a
categorical generator and inherit much of the structural vocabulary of
\textcite{song2021score} despite working in a different state space.

\section{Predictor--corrector samplers}
\label{lit:sec:pred-corr}

\textcite{song2021score} introduce predictor--corrector samplers as a
practical way to reduce the discretisation error of reverse-SDE
integrators. A predictor step advances the chain from time $t$ to time
$t-\Delta t$ using a numerical scheme for the reverse SDE or ODE. A
corrector step then runs one or more iterations of an MCMC kernel that
targets $q_{t-\Delta t}$, holding the time index fixed. The canonical
corrector is overdamped Langevin dynamics:
\[
  x \;\leftarrow\; x + \tfrac{\eta_t^2}{2}\,\stheta(x,t-\Delta t) + \eta_t\,z,
  \qquad z\sim\mathcal{N}(0,I),
\]
applied for $K$ iterations with a step size $\eta_t$ scaled to the local
noise level. \textcite{song2021score} report that, on image benchmarks,
replacing pure-predictor sampling with predictor--corrector sampling at
fixed total NFE consistently lowers Fr\'echet Inception Distance.

The conceptual structure is general and is the immediate ancestor of the
masked-diffusion correctors discussed in
Chapter~\ref{lit:ch:correctors}. Three properties are worth isolating.
First, the predictor and corrector consume the same trained network; they
differ only in how that network's output is used. Second, the corrector
targets a fixed marginal at the current time, so it does not violate the
time decomposition supplied by the forward process. Third, both steps cost
one or more model evaluations, so allocating effort between predictor and
corrector trades against allocating effort to additional predictor
steps. The last property already foreshadows fixed-budget timing: with a
finite NFE budget, where should corrector calls be placed along the
trajectory?

In the continuous case this question is partly answered by the Langevin
diffusion's local mixing rate, which in turn is governed by the conditioning
of $q_t$. At noise levels where $q_t$ is well concentrated and isotropic, a
Langevin corrector mixes quickly and a few iterations suffice; at
intermediate noise levels, where the marginal is multi-modal, mixing is
slower and additional corrector steps may be more valuable. The thesis
question asks for the analogue of this allocation principle in the discrete
masked setting, where the corrector is no longer a Langevin step but a fixed
informed Markov kernel on tokens.

\section{Deterministic and accelerated samplers}
\label{lit:sec:fast-samplers}

Because each predictor or corrector step is a model evaluation, the dominant
inference cost in diffusion models is the number of function evaluations
(NFE). Two families of accelerated samplers reduce NFE without retraining
the score network.

The first uses non-Markovian deterministic trajectories.
\textcite{song2020ddim} introduce DDIM, a family of samplers that share
marginals with the DDPM forward process but allow deterministic transitions
and consistent latent--sample correspondences. With $\eta=0$ the DDIM
sampler is invertible and acts as a pseudo-encoder; with $\eta=1$ it
recovers DDPM. DDIM trades stochasticity for fewer steps and is widely used
as a low-NFE baseline.

The second family exploits the semi-linear structure of the diffusion ODE.
\textcite{lu2022dpmsolver} introduce DPM-Solver, an exponential integrator
tailored to the structure
$\dif x_t/\dif t = f(t)\,x_t + g(t)\,\stheta(x_t,t)$ with affine $f,g$;
the resulting samplers achieve high-quality samples in roughly ten function
evaluations on image benchmarks. Subsequent work develops higher-order
solvers and adaptive step-size control, and the related literature on
solver-aware schedule learning tunes the discretisation grid for a given
solver and dataset.

What these accelerated samplers share is that they treat the trained score
network as a black box and recompose the inference trajectory: where to
sample time points, how to integrate, and how to allocate function
evaluations. This is the same family of design questions that the present
thesis instantiates for masked-diffusion correctors. The continuous
literature has already established that, with a fixed score network,
inference-time scheduling is a meaningful research object; the discrete
literature considered in subsequent chapters has begun to develop its own
counterpart.

\section{Inference-time guidance}
\label{lit:sec:guidance}

A third category of inference-time intervention modifies the sampling
distribution rather than the schedule. \textcite{dhariwal2021diffusionbeats}
use an auxiliary classifier $p_\phi(y\mid x_t)$ trained on noisy inputs and
bias the score during reverse sampling toward class $y$:
\[
  \tilde s(x_t,t,y)
  = \stheta(x_t,t) + w\,\nabla_{x_t}\log p_\phi(y\mid x_t),
\]
with a guidance scale $w\geq 0$. \textcite{ho2022cfg} show that the
auxiliary classifier can be removed by training the score network on both
conditional and unconditional inputs (with random label dropout) and
combining at sampling:
\[
  \tilde s(x_t,t,y)
  = (1+w)\,\stheta(x_t,t,y) - w\,\stheta(x_t,t,\varnothing),
\]
where $\varnothing$ is a null condition. Classifier-free guidance is now the
default mechanism for conditional diffusion across modalities.

Guidance is interesting for the present thesis because it is structurally
similar to corrector ideas: both modify the sampler at inference without
retraining the base model, and both can be tuned per noise level. Unlike
correctors, however, guidance changes the marginal distribution itself
rather than reducing approximation error toward an unchanged marginal. The
analogue in masked diffusion is reward-guided refinement
(Chapter~\ref{lit:ch:correctors}, \S\ref{lit:sec:prism}), which similarly
biases the inference distribution toward higher-reward outputs.

\section{Why continuous diffusion theory does not transfer directly to language}
\label{lit:sec:to-discrete}

The continuous theory is mathematically convenient because the state space
is Euclidean. Small perturbations are meaningful, gradients exist, the
score $\nabla_x\log q_t$ is well defined, and Gaussian transition kernels
are tractable. Text is different. A token sequence lives in the finite
product space $\mathcal{V}^L$, where $|\mathcal{V}|$ is a finite vocabulary
size and $L$ is sequence length; for a typical language model
$|\mathcal{V}|$ is in the tens of thousands and $L$ in the thousands. There
is no infinitesimal move from one token to another, no canonical metric
tensor on $\mathcal{V}^L$ that would make `local Gaussian noise' meaningful,
and no continuous score field whose gradient one could integrate.

Three options have been pursued. The first embeds tokens in a continuous
space, runs a continuous diffusion in the embedding, and decodes back to
tokens at sampling time. Embedding-space approaches introduce an
additional decoding problem and decouple the diffusion process from the
discrete data distribution; the remainder of this review focuses on
diffusion processes defined directly on discrete token states. The
second replaces the continuous SDE with a continuous-time Markov chain
on $\mathcal{V}^L$, with rate matrices governing categorical transitions.
The third uses a discrete-time categorical Markov chain. The latter two
preserve the discrete data distribution at every step.

Chapter~\ref{lit:ch:discrete-diffusion} reviews the discrete-time and
continuous-time discrete diffusion literature.
Chapter~\ref{lit:ch:masked-diffusion} restricts attention to the masked or
absorbing-state special case, the empirical setting for the rest of the
thesis. The continuous theory remains useful, but principally as a source of
structural vocabulary---predictor versus corrector, schedule as design
variable, NFE as the operational cost---rather than as an importable theorem
stack.


%% =====================================================================
%% Chapter 3 -- Discrete Diffusion Models
%% =====================================================================

\chapter{Discrete Diffusion Models}
\label{lit:ch:discrete-diffusion}

The discrete state space changes the type of diffusion process available
and, consequently, the type of inference-time intervention that is
meaningful. This chapter reviews the discrete diffusion literature with a
focus on three questions: how the forward corruption is constructed on a
finite alphabet; how the reverse process is parameterised and trained; and
how the unavoidable factorisation of the reverse over positions interacts
with sampling error. The factorisation question motivates the move to
inference-time correction in Chapter~\ref{lit:ch:correctors}.

\section{Categorical state spaces and discrete corruption}
\label{lit:sec:categorical}

Let $\calV$ be a finite alphabet of size $|\calV|$ and let
$\mathcal{X}=\calV^{\seqlen}$ be the product space of length-$\seqlen$
sequences. The data distribution $\ptrue$ is a probability mass function on
$\mathcal{X}$. Unlike $\mathbb{R}^d$, the space $\mathcal{X}$ has no group
structure compatible with `addition of small noise': two tokens are either
equal or unequal, and there is no canonical interpolation between them.

The standard discrete diffusion construction therefore replaces additive
Gaussian noise by a categorical Markov chain on $\calV$ (or, in some
constructions, on an enlarged state space such as $\calVm = \calV\cup\{\Mask\}$).
A forward chain at time $t$ is given by a transition matrix
$Q_t\in\mathbb{R}^{|\calV|\times|\calV|}$ whose rows are probability
distributions over the next state. The chain factorises across positions:
\[
  \qfwd(\zt \mid \clean)
  = \prod_{\ell=1}^{\seqlen} \qfwd\big(\zt^\ell \mid \clean^\ell\big),
\]
where $\qfwd(\zt^\ell \mid \clean^\ell)$ is the row of the cumulative
transition matrix $\bar Q_t = Q_1 Q_2 \cdots Q_t$ indexed by $\clean^\ell$.
The position-wise factorisation is a property of the forward process, not of
the data, and is the structural reason that discrete diffusion can be
trained efficiently: one can sample $\zt$ given $\clean$ in one step
regardless of the inter-position dependencies of $\ptrue$.

\section{Discrete-time and continuous-time formulations}
\label{lit:sec:dt-ct}

There are two common parameterisations of the forward chain. The
discrete-time formulation chooses a sequence of stochastic matrices
$Q_1,\ldots,Q_T$ with $\bar Q_t = \prod_{s\leq t}Q_s$ and views generation
as a backward chain over the index set $\{T,T-1,\ldots,1,0\}$. The
continuous-time formulation specifies a rate matrix
$R_t\in\mathbb{R}^{|\calV|\times|\calV|}$ with non-negative off-diagonals
and zero row sums; the marginal $\bar Q_t$ then satisfies the Kolmogorov
forward equation
\[
  \od{\bar Q_t}{t} = \bar Q_t R_t, \qquad \bar Q_0 = I.
\]
The continuous-time view is closer in spirit to the SDE formulation of
\S\ref{lit:sec:sde} and is the framework in which several recent papers
state results.

The reverse process takes the sampled $z_T \sim \qfwd(\cdot)$ and traces
back toward $z_0$ using a learned conditional distribution
$\ptheta(z_{t-1}\mid \zt)$ at each step. The exact reverse conditional under
the forward chain is
\[
  \qfwd(z_{t-1}\mid \zt, \clean)
  = \frac{\qfwd(\zt\mid z_{t-1})\,\qfwd(z_{t-1}\mid \clean)}{\qfwd(\zt\mid \clean)},
\]
and depends on the unobserved $\clean$. In practice the reverse model is
trained either to predict $\clean$ from $\zt$ (an $x_0$-prediction
parameterisation) or to estimate position-wise transition ratios. Both
parameterisations factorise the reverse over positions, which is the source
of the inference-time approximation gap discussed in
\S\ref{lit:sec:factorization} and \S\ref{lit:sec:lz}.

\section{D3PM and structured discrete diffusion}
\label{lit:sec:d3pm}

\textcite{austin2021d3pm} introduce structured denoising diffusion
probabilistic models (D3PMs) and provide the first systematic treatment of
discrete diffusion across families of forward kernels. They consider four
families of $Q_t$ matrices. \emph{Uniform} kernels mix smoothly toward the
uniform distribution on $\calV$. \emph{Gaussian-like} kernels use transition
probabilities induced by a discretised Gaussian on an ordered or
otherwise structured discrete state space, producing locality structure
on the alphabet. \emph{Absorbing} kernels
allow each token to transition to a special $\Mask$ state, making the chain
monotone in a masked set. \emph{Discretised continuous-state} kernels mimic
a continuous diffusion through binning. Each family produces a different
inductive bias and a different decoding behaviour.

The reverse-process loss in \textcite{austin2021d3pm} is a hybrid of an
evidence lower bound and an auxiliary cross-entropy on the predicted $x$
from each $\zt$. The auxiliary loss empirically improves training stability
and provides a direct supervision signal at every time step. The D3PM
construction makes explicit the position-wise factorisation of both the
forward chain and the reverse model, which simplifies both training and
sampling and is inherited by all subsequent discrete diffusion families.

The absorbing kernel introduced by \textcite{austin2021d3pm} is the
immediate ancestor of masked diffusion language models. Because absorbing
transitions are monotone---a token, once moved to $\Mask$, remains masked
under further noising---the chain has a particularly simple structure: the
set of masked positions at time $t$ is a superset of the set at time
$t' < t$. This monotonicity is what makes masked diffusion empirically
attractive for language and is treated in detail in
Chapter~\ref{lit:ch:masked-diffusion}.

\section{SEDD and the score-entropy view}
\label{lit:sec:sedd}

\textcite{lou2024sedd} reformulate discrete diffusion through ratio
estimation. For a continuous-time Markov chain with rate matrix $R$, the
reverse generator depends on the ratios $\qfwd(z')/\qfwd(z)$ for pairs of
adjacent states, and \textcite{lou2024sedd} introduce a parameterisation
that estimates these ratios directly. They define the score-entropy loss,
a discrete analogue of denoising score matching, which is consistent under
mild assumptions and admits closed-form gradients for several common rate
matrices.

Two structural advantages follow. First, ratio estimation avoids the need
for an explicit normalisation constant in the discrete reverse
parameterisation. Second, the score-entropy loss is convex in the predicted
ratios, in contrast to the evidence-lower-bound objective which is
non-convex even at the function-class level. \textcite{lou2024sedd} report
empirical gains on language modelling benchmarks, although later models
with absorbing-state forward processes
(Chapter~\ref{lit:ch:masked-diffusion}) report lower perplexity at
comparable parameter counts in the benchmark suites used by subsequent
masked-diffusion papers.

The SEDD parameterisation is conceptually closest to the continuous
score-based view of \S\ref{lit:sec:score-matching}: both estimate a
non-normalised quantity associated with the forward process and use it to
drive the reverse dynamics. The main difference is that the discrete
`score' is a finite collection of ratios rather than a vector field. SEDD
also illustrates that the choice of training loss for discrete diffusion is
genuinely open: D3PM uses an ELBO with auxiliary cross-entropy, SEDD uses
score-entropy, and masked diffusion language models
(Chapter~\ref{lit:ch:masked-diffusion}) use a tightened cross-entropy
bound.

\section{Discrete flow matching}
\label{lit:sec:dfm}

\textcite{gat2024dfm} extend the flow matching framework of continuous
diffusion to discrete states. Flow matching parameterises a probability path
$\pi_t$ between a source distribution $\pi_0$ (typically a simple reference)
and a target distribution $\pi_1=\ptrue$, and trains a network to match the
conditional vector field along this path. In the discrete version, the path
is a sequence of categorical distributions and the conditional vector field
is replaced by a generator of a categorical jump process.

Two consequences are relevant for the present thesis. First, discrete flow
matching is the most general framework currently available for
parameterising discrete diffusion: D3PM corresponds to specific source and
path choices, masked diffusion to absorbing source and product path. Second,
the framework makes explicit that the source distribution is a design
variable. In particular, masked diffusion's choice of source---the all-mask
sequence---is one of many. \textcite{gat2024dfm} also study uniform and
other categorical sources, and the empirical advantages of the absorbing
source are not automatic from the framework but emerge from the interaction
of source choice, path geometry, and parameterisation.

\section{Factorised reverse approximations and inference error}
\label{lit:sec:factorization}

A constraint shared by all discrete diffusion models considered in this
chapter is that the reverse model factorises over positions:
\[
  \pfact(z_{t-1}\mid \zt)
  = \prod_{\ell=1}^{\seqlen}\ptheta(z_{t-1}^\ell\mid \zt).
\]
This is computationally necessary: the joint conditional distribution over
$\seqlen$ positions has $|\calV|^{\seqlen}$ atoms, and any practical decoder
must factorise. But the true reverse conditional $\qfwd(z_{t-1}\mid \zt)$ is
generally not factorised across positions, even when the forward chain is:
marginalising over the unknown $\clean$ introduces dependence between
positions, and that dependence is exactly the data dependence the reverse
should model.

This factorisation gap is the discrete-diffusion analogue of the
discretisation error that predictor--corrector samplers were designed to
address in \S\ref{lit:sec:pred-corr}. In the continuous-time limit each
step acts on at most one position and the factorised parameterisation
aligns with the ideal reverse generator. At any finite horizon used in
practice, the factorised reverse incurs a bias relative to the joint
true reverse conditional that the trained network cannot absorb
pointwise. Inference-time mechanisms that operate jointly across
positions, or that refine one position at a time after others have been
committed, are the natural targets for reducing this finite-$T$ bias.

The observation places discrete diffusion in a different inference-design
regime from its continuous counterpart. A continuous-time SDE can in
principle be integrated to within numerical precision, and a Langevin
corrector is then an error-reducing tool against discretisation alone. A
practical discrete-time factorised reverse, run at the horizons used in
practice, has a non-trivial gap between the learned per-position
predictions and the joint reverse conditional, and this gap motivates
joint refinement during inference rather than further time refinement
alone.

\section{The Lavenant--Zanella error decomposition}
\label{lit:sec:lz}

\textcite{lavenant2025lz} make this gap precise for masked diffusion language
models with factorised reverses. They establish a quantitative error bound
whose two leading terms admit a clean operational interpretation: a
\emph{learning} error $\Elearn$ that captures the gap between the trained
denoiser $\ptheta(x\mid \zt)$ and the Bayes-optimal denoiser
$q(x\mid \zt)$, and a \emph{factorisation} error $\Efact$ that captures
the gap between the factorised reverse $\pfact$ and the true reverse
$\qfwd(z_{t-1}\mid \zt)$ even when the denoiser is exact.

The decomposition has two implications relevant to the present thesis.
First, $\Elearn$ scales with the quality of the trained denoiser and is a
property of training, not of inference. Second, $\Efact$ is structurally
inference-time: it depends on the time discretisation, the sampling
procedure, and any inference-time intervention applied to the trajectory.
The decomposition itself is structural; whether either term dominates in
a given model is an empirical question outside the scope of the bound.
The separation suggests that inference-time mechanisms should be analysed
separately from denoiser training error, even when the two cannot always
be cleanly disentangled in practice.

\textcite{lavenant2025lz} additionally study optimal time schedules for
masked diffusion under their bound and derive prescriptions for the rate at
which mask probability should change with $t$. The optimal schedules are
not uniform-in-time, and the deviation from uniform depends on the data
distribution and the noise level via specific functional dependencies. The
schedule-optimisation result is conceptually important for the thesis:
even at the level of the predictor schedule, the choice of where to spend
inference effort is not benign. Corrector timing, the object of the present
thesis, is the analogous question one rung higher---with the predictor
schedule fixed, where to spend a separate corrector budget.

\textcite{chen2025optimalschedules} subsequently strengthen the analysis.
They derive a tight characterisation of the divergence between the true and
the sampled distribution under masked diffusion as a function of the
unmasking schedule, expose a connection to univariate function
approximation theory, and prove both lower and upper bounds. They also show
that, without strong a priori knowledge of the data distribution, no
schedule can compete with the optimal schedule on every distribution,
while in benign settings characterised by total correlation and dual total
correlation $\mathcal{O}(\log L)$ steps suffice. The result is the
strongest current rigorous statement of why scheduling matters for masked
diffusion sampling, and it bounds what a fixed predictor schedule can
achieve on its own.

\section{The case for inference-time correction}
\label{lit:sec:case-correction}

Three observations from the foregoing are worth assembling. First, the
practical discrete diffusion models reviewed here rely on tractable local
or factorised reverse parameterisations. Second, the resulting
finite-horizon sampling error is structurally inference-time, separate
from training error in the sense that it persists when the trained
denoiser is held fixed. Third, the recent error-decomposition and
schedule-theoretic results
\parencite{lavenant2025lz,chen2025optimalschedules} expose this gap and
tie it to specific properties of the inference procedure.

These observations motivate the move to inference-time correction. A
corrector is a fixed Markov transition that operates on the current state
and is intended to reduce finite-horizon sampling error at the cost of
additional NFE. Whether
such a corrector exists for a given discrete diffusion family, what its
action set is, and how to schedule its use are separate questions and form
the subject of Chapter~\ref{lit:ch:correctors}. Before that,
Chapter~\ref{lit:ch:masked-diffusion} narrows attention to the masked or
absorbing-state special case, which is the empirical setting for the
remainder of the thesis.


%% =====================================================================
%% Chapter 4 -- Masked Diffusion Language Models
%% =====================================================================

\chapter{Masked Diffusion Language Models}
\label{lit:ch:masked-diffusion}

Masked diffusion language models (MDLMs) are the absorbing-state special
case of discrete diffusion. The forward process either keeps each token as
is or replaces it with a single dedicated $\Mask$ symbol, never returning a
previously masked position to the alphabet, and the reverse process learns
to fill in masked positions from the surrounding context. This chapter
develops the masked-diffusion construction, examines where its monotone
absorbing structure simplifies training and sampling relative to alternative
discrete diffusion families, surveys the recent scaling efforts toward
language model regimes, and introduces the inference-time mechanics---token
ordering, unmasking schedules, committed-token revisability---that will be
needed for the remasking and corrector literature in
Chapter~\ref{lit:ch:correctors}.

\section{The absorbing-mask process}
\label{lit:sec:absorbing}

Let $\calV$ be the alphabet, $\seqlen$ the sequence length, and $\Mask$ a
distinguished symbol not in $\calV$. The masked-diffusion forward chain
operates on the extended state space $\calVm = \calV\cup\{\Mask\}$. Fix a
non-increasing schedule $\alpha_t : [0,1]\to[0,1]$ with $\alpha_0 = 1$ and
$\alpha_1 = 0$, interpreted as the fraction of positions that remain
unmasked at time $t$. The forward process replaces each clean token by
$\Mask$ independently with probability $1-\alpha_t$:
\[
  \qfwd\big(\zt^\ell \mid \clean^\ell\big)
  = \alpha_t\,\delta_{\clean^\ell} + (1-\alpha_t)\,\delta_{\Mask},
  \qquad \ell=1,\ldots,\seqlen.
\]
The chain is absorbing: once a position is marked $\Mask$, it remains so
under further noising. The marginal $\qfwd(\zt\mid \clean)$ therefore
depends only on the random subset $M_t \subseteq \{1,\ldots,\seqlen\}$ of
currently masked positions, and the position-wise transitions are
independent.

The reverse model is conventionally parameterised as an $x$-predictor:
given a partially masked sequence $\zt$, the network outputs a categorical
distribution $\ptheta(\cdot\mid \zt)$ over $\calV^{\seqlen}$ that factorises
across positions and is supported only on the masked positions (committed
tokens are passed through). Sampling at time $t$ then unmasks a
designer-chosen subset of masked positions by drawing from the predicted
marginal at each. The schedule $\alpha_t$ controls the expected fraction of
masked positions at time $t$ but does not specify which positions are
unmasked at any given step; the latter is the unmasking policy of
\S\ref{lit:sec:token-ordering}.

The forward construction is mathematically simple and computationally
cheap: a single sample of $(\clean,t,U)$ with
$U\sim\mathrm{Uniform}[0,1]^{\seqlen}$ produces $\zt$ in closed form by
thresholding $U$ against $\alpha_t$ position-wise. This shortcut makes
large-scale MDLM training feasible and is the masked-diffusion analogue of
the DDPM Gaussian shortcut from \S\ref{lit:sec:fwd-rev}.

\section{MDLM and the simplified objective}
\label{lit:sec:mdlm}

\textcite{sahoo2024mdlm} establish a tight evidence lower bound for masked
diffusion that admits a particularly simple training loss. Writing the
position-wise cross-entropy of the predicted denoiser against the
ground-truth token at position $\ell$ as
$\mathrm{CE}\!\big(\clean^\ell, \ptheta(\cdot\mid \zt)\big)$, they show that
the negative ELBO collapses to a weighted sum
\[
  -\mathrm{ELBO}
  \;\leq\;
  \E_{\clean,t,\zt}\,
  \sum_{\ell:\, \zt^\ell = \Mask}
    w_t\,\mathrm{CE}\!\big(\clean^\ell, \ptheta(\cdot\mid \zt)\big),
\]
where $w_t = -\dot\alpha_t/(1-\alpha_t)$ is a closed-form weight depending
only on the schedule and where $\dot\alpha_t$ is the time derivative of the
schedule. The unmasked positions contribute zero loss because the forward
chain leaves them invariant. The training procedure is therefore a weighted
masked-language-model regression with weights dictated by the schedule.
This is the key empirical advantage that \textcite{sahoo2024mdlm} document:
a simple training objective with a tight ELBO, requiring no auxiliary
cross-entropy term and no complex sampling-time machinery.

\textcite{sahoo2024mdlm} also derive a parameterisation of the reverse
transition whose coefficients are fixed by the absorbing schedule rather
than left as tunable hyperparameters. Sampling reduces to a straightforward
iteration: at each step the predicted distribution at masked positions is
sampled, the masked set is updated, and the schedule advances.

A practical consequence of the MDLM construction is that the trained
denoiser is a single neural network that maps any partially masked sequence
$\zt$ to per-position distributions over $\calV$. Once trained, that
network supports a family of inference strategies: vanilla iterative
unmasking, batched parallel decoding, conditional infilling, and---relevantly
for this thesis---remasking and corrector samplers
(Chapter~\ref{lit:ch:correctors}) that revisit committed positions.

\section{Scoped advantages of masked diffusion over alternative discrete
families}
\label{lit:sec:masked-advantages}

Among recent language-focused discrete diffusion papers, absorbing-mask
models are the most prominent high-performing family at present, but the
literature does not establish universal dominance. The differences
between masked and non-masked discrete diffusion are best understood as a
small set of structural advantages and trade-offs.

The clearest advantage is that the absorbing forward chain is monotone in
the masked set, which simplifies the ELBO. Under a uniform forward chain
(D3PM-uniform), every position can transition to any token at every step,
and the exact reverse conditional has support on the full
$|\calV|^{\seqlen}$ space. The factorised reverse therefore has to fit a
more entangled marginal. Under absorbing forward chains, the masked
positions are clearly identified and the reverse only needs to fit the data
marginal at masked positions conditional on the committed tokens, which is
closer to the masked-language-model setup that pretraining literature has
converged on.

\textcite{lou2024sedd} report that the score-entropy training objective
with absorbing-mask forward chains attains the best perplexity in their
suite, but the gap between absorbing and uniform varies across model scales
and is sensitive to training compute. \textcite{austin2021d3pm} note that
the absorbing kernel benefits from a compatibility between forward chain
and supervision: the network is trained to predict $\clean$ given $\zt$,
and the absorbing chain ensures that, at any $t$, the supervision is
well-defined at all masked positions and trivially correct at unmasked
positions.

A second advantage is that masked-diffusion sampling never has to undo a
wrong substitution. Under a non-absorbing chain such as D3PM-uniform, a
position can be assigned a token, then re-corrupted, then assigned a
different token; the reverse process must implicitly handle this revision.
Under absorbing chains, once a position is unmasked, the standard sampler
leaves it alone---unless an explicit revisability mechanism is added. This
simplification of the trajectory geometry is why masked diffusion is a
particularly clean substrate for studying additional inference-time
corrections.

There are also limitations. Masked diffusion gives up the within-step token
revision that uniform discrete diffusion supports `for free'. Adding
revisability back---through remasking samplers, self-correction backends,
or informed correctors---is the topic of Chapter~\ref{lit:ch:correctors}.
Masked diffusion also commits, by construction, to a particular path
through state space (positions become unmasked monotonically), which
constrains the reverse trajectories that can be expressed; \textcite{wu2026dsl}
and \textcite{huang2025remedi} discuss this constraint explicitly and
propose remasking-aware training to relax it.

The current literature consensus is therefore best summarised as follows.
Masked diffusion offers a particularly tight ELBO, a simple training
objective, and a clean trajectory geometry. It does not, on the available
evidence, dominate alternative discrete diffusion formulations on every
axis, and its monotone-mask constraint introduces specific limitations that
subsequent inference-time mechanisms are designed to address.

\section{Scaling masked diffusion to language model regimes}
\label{lit:sec:scaling}

Several recent works pretrain masked diffusion language models at scales
comparable to mid-sized autoregressive language models.
\textcite{nie2025llada} introduce LLaDA, an 8B-parameter masked diffusion
language model trained on 2.3T tokens with a curriculum that mixes
mask-prediction and next-token-prediction objectives. The authors report
parity with LLaMA-3 8B-Instruct on several benchmarks and document
diffusion-specific inference properties at this scale: bidirectional
context use, native infilling, and parallel decoding.

\textcite{ye2025dream} introduce Dream-7B, a 7B masked diffusion language
model with an alternative pretraining recipe. The paper reports parity
with autoregressive baselines on selected reasoning and coding benchmarks
while supporting native parallel and infilling modes.
\textcite{sahoo2025eso} introduce Esoteric Language Models (Eso-LMs),
which interpolate between autoregressive and masked-diffusion regimes
during training so that the resulting model supports both modes at
inference. The claimed deployment tradeoff is that a single Eso-LM can be
configured per-input as autoregressive (for KV-cached low-latency
single-user generation) or as masked-diffusion (for parallel decoding
and infilling).

These scaling efforts establish that the masked-diffusion training recipe
is compatible with the same scale of compute as autoregressive language
modelling and that the resulting models inherit the inference-time
flexibility of masked diffusion. They do not yet establish parity at the
largest scales: the published frontier-scale autoregressive models remain
substantially larger than published MDLMs at the time of writing, and
direct comparisons at matched compute are limited. The relevance for the
present thesis is that the empirical backend used in the contribution
chapters (an MDLM-style ProSeCo OpenWebText model) inherits the
inference-time mechanics---monotone mask, factorised reverse, schedule-based
sampling---unchanged from this scaling literature.

\section{Token ordering and unmasking schedules}
\label{lit:sec:token-ordering}

The masked-diffusion sampler must specify, at each step, both how many
positions to unmask and which positions to unmask. The first is governed by
the schedule $\alpha_t$ and is essentially a property of the forward chain;
the second is the unmasking policy and is an inference-time degree of
freedom not fixed by training.

\textcite{kim2025tokenordering} study the sensitivity of MDLM sample
quality to the unmasking policy. They show empirically that the canonical
training-time policy---random masking---is not optimal at inference, and
that a planning-time policy that prefers high-confidence tokens improves
sample quality. Their analysis decomposes the error into a `train for the
worst, plan for the best' principle: train against worst-case mask sets to
ensure robustness, then exploit confidence-aware unmasking at inference.
Subsequent work refines this picture. \textcite{kim2025klass} introduce
KL-guided unmasking, in which the policy minimises KL divergence to a
reference distribution; \textcite{hong2025upo} train the unmasking policy
directly to optimise downstream reward; \textcite{he2025mdpo} formalise the
discrepancy between training and inference unmasking distributions and
propose a training-time correction.

\textcite{benhamu2025ebsampler} introduce entropy-bounded sampling, in
which the number of positions unmasked per step is determined by a target
entropy on the predicted token distribution. \textcite{chen2025denoisingentropy}
study uncertainty-driven decoding paths and show that quantifying
uncertainty per position can improve sample quality.
\textcite{luxembourg2025plan} propose dilated unmasking schedules that
target inference speed by parallelising appropriately conditioned positions.
These methods all share a common principle: the predictor schedule and the
unmasking policy together are tunable inference-time machinery, separate
from the trained denoiser.

The thesis problem is one rung higher than this token-ordering literature.
The token-ordering literature asks `at each unmasking step, which positions
should be unmasked?' with a fixed predictor budget. The corrector-timing
problem asks `at which time steps should an additional refinement
transition be applied?' with a fixed predictor schedule, a fixed corrector
kernel, and a fixed budget for corrector calls. Token ordering and corrector
timing are complementary---one chooses how to unmask, the other chooses
when to refine---and both are inference-time scheduling questions over a
fixed denoiser.

\textcite{chen2025optimalschedules}, in addition to their schedule-theoretic
results discussed in \S\ref{lit:sec:lz}, give the cleanest current statement
of why these schedule questions are meaningful: even with a Bayes-optimal
denoiser, the choice of unmasking schedule affects the divergence between
the sampled and true distributions, and this dependence cannot be removed
by training alone. Inference-time scheduling is therefore not a heuristic
on top of a finished model; it is a structural component of the generative
procedure.

\section{The clean trajectory perspective}
\label{lit:sec:clean-trajectory}

Masked diffusion gives a particularly clean trajectory on which to operate.
A sample path is a sequence
\[
  z_T,\, z_{T-1},\, \ldots,\, z_0,
\]
with $z_T$ either fully masked or near-fully masked and $z_0$ a completed
token sequence. The state at step $t$ is fully determined by the masked
set $M_t$ and the committed tokens at positions outside $M_t$, with
$M_T \supseteq M_{T-1} \supseteq \cdots \supseteq M_0$ in the absence of
remasking. Inference-time interventions---accelerated solvers, alternative
unmasking policies, remasking, correctors, search procedures---all operate
on this trajectory.

The cleanness has three operational implications for the rest of the
thesis. First, the corrector kernel has a well-defined action set at each
step, namely a subset of positions either still masked or, in
self-correcting backends, eligible for revision. Second, intermediate states
$\zt$ are explicit: any signal computed at step $t$ has a precise referent,
in contrast to continuous diffusion where the intermediate state is a
sample from a marginal that the network can only estimate. Third, the cost
of inserting a corrector call at step $t$ is well-defined: it is the cost
of one or more additional forward passes, conditional on $\zt$. The latter
two properties are what make a paired schedule-gain definition
$G(S) = F(z_0^{S}) - F(z_0^{\varnothing})$ a meaningful object even before
a particular corrector or quality functional is fixed.

The next chapter introduces the inference-time mechanisms that operate on
this trajectory: remasking samplers, uncertainty-driven decoding, informed
correctors, self-correcting backends, quality-guided refinement, and
trajectory search. The chapter ends by motivating the timing question that
the contribution chapters formalise.


%% =====================================================================
%% Chapter 5 -- Remasking, Self-Correction, and Informed Correctors
%% =====================================================================

\chapter{Remasking, Self-Correction, and Informed Correctors}
\label{lit:ch:correctors}

The previous chapters established that masked diffusion language models
leave a structural gap between the factorised reverse process used at
inference and the true reverse conditional, and that this gap is not closed
by training. Inference-time mechanisms that operate on the masked
trajectory are designed to reduce this gap, and the recent literature has
produced a substantial number of them. This chapter organises that
literature into a taxonomy useful for the present thesis---remasking
samplers, uncertainty-driven decoding, informed correctors, self-correction
backends, quality-guided refinement, and trajectory search---and then
narrows attention to the informed-corrector lineage that is the immediate
ancestor of corrector timing. A short concluding section sketches the
structural reasons why corrector timing is itself a research question,
deferring the formal framework to Chapter~\ref{ch:contribution}.

\section{A taxonomy of inference-time mechanisms}
\label{lit:sec:taxonomy}

The methods reviewed in this chapter all modify inference-time behaviour,
but at different layers. It is useful to distinguish them up front. A
\emph{base masked diffusion language model} supplies the trained denoiser
and a default sampling procedure (\S\ref{lit:sec:mdlm}). A \emph{remasking
sampler} changes the sampling procedure: positions previously committed
may be returned to the masked state and resampled, with the net effect
being a non-monotone trajectory through state space. An \emph{informed
corrector} is a fixed Markov transition that operates on the current state
at chosen times, refining a designated action set without changing the
underlying sampler. A \emph{self-correction backend} is a base model
jointly trained to support both unmasking and refinement, so that the
corrector and the predictor share parameters. \emph{Quality-guided
refinement} biases the inference distribution toward higher-reward outputs
using learned per-token quality scores. \emph{Trajectory search}, finally,
branches over candidate trajectories or candidates at each step and uses
external feedback (a reward, a verifier, or true schedule-gain information)
to keep the best.

These categories overlap. A self-correction backend supplies a specific
informed corrector (the trained refinement head) and changes the action set
(committed tokens become eligible for revision). A remasking sampler can
be reinterpreted as a designed corrector with a particular action set
(positions to unmask again). Quality-guided refinement can supply a
per-step signal that an informed corrector or a remasking sampler can use
to choose what to refine. The thesis question---corrector timing---takes
one specific configuration: a fixed predictor schedule, a fixed informed
corrector, and a fixed corrector-placement budget, and asks at which time
steps the corrector should be applied.

This isolation matters because the literature does not always distinguish
the layers cleanly. The taxonomy in Table~\ref{lit:tab:taxonomy} is the
working vocabulary used for the rest of the chapter.

\begin{table}[t]
\centering
\small
\begin{tabular}{p{3.6cm} p{6.2cm} p{4.6cm}}
\toprule
Layer & What changes at inference & Examples \\
\midrule
Base model
  & Sampler and denoiser fixed
  & MDLM, SEDD, LLaDA, Dream \\
Remasking sampler
  & Trajectory itself; positions can be re-masked
  & ReMDM, RemeDi, DSL \\
Self-correction backend
  & Sampler refines committed tokens jointly with unmasking
  & ProSeCo \\
Informed corrector
  & Fixed Markov kernel applied at chosen steps to a single trajectory
  & Zhao et al.\ correctors \\
Quality-guided refinement
  & Inference biased by learned quality scores
  & PRISM, IterRef, S$^3$ \\
Trajectory search / TTS
  & Branch over candidates or particles with external feedback
  & IterRef, S$^3$, PG-DLM, Feynman--Kac correctors \\
\bottomrule
\end{tabular}
\caption{Working taxonomy of inference-time mechanisms for masked diffusion
language models. Layers are distinct in principle and overlap in practice;
the corrector-timing question of this thesis fixes the base model, the
corrector kernel, and the corrector budget, and varies only the placement
schedule.}
\label{lit:tab:taxonomy}
\end{table}

\section{Remasking samplers}
\label{lit:sec:remasking}

Remasking samplers extend masked diffusion by allowing tokens to be
returned to the masked state mid-trajectory. \textcite{wang2025remdm}
introduce ReMDM, the canonical remasking-sampler formulation. ReMDM
specifies a probability schedule for re-masking each currently committed
position at each step, parameterised by a per-step rate that is itself an
inference-time hyperparameter. The denoiser is the standard MDLM-trained
network; the sampler does the work. \textcite{wang2025remdm} show that a
calibrated re-masking schedule improves sample quality at the cost of
additional sampling steps, and that the technique benefits from being
combined with confidence-based unmasking policies.

\textcite{huang2025remedi} introduce RemeDi (Remasking-enabled Diffusion
Language Model), in which the denoiser jointly predicts both token
distributions and per-token confidence scores, and the sampler uses the
confidence scores to remask uncertain tokens at each step. The training
objective is augmented with a confidence head that is supervised against
held-out signals. The mechanism is a remasking sampler with an integrated
training-time signal, sitting between the pure-sampler ReMDM and the
joint-training ProSeCo (\S\ref{lit:sec:proseco}).

\textcite{wu2026dsl} propose Discrete Stochastic Localization (DSL), which
trains a single SNR-invariant denoiser across a continuum of corruption
levels including the remasking-induced `going back to lower SNR'
transitions. The argument is that current MDLM training does not match the
input distribution seen by remasking samplers, and that bridging this
distributional gap during training improves sample quality at moderate
compute cost. \textcite{zhang2026cdlm} introduce CDLM (Corrective Diffusion
Language Models), which observes that MDLMs often fail to assign lower
confidence to incorrect tokens after they have been committed, and propose
a training-time augmentation that explicitly supervises corrective
behaviour. The resulting model produces calibrated confidence at committed
positions, which any remasking or corrector mechanism can then use as a
signal.

The boundary between remasking samplers and self-correction backends
(\S\ref{lit:sec:proseco}) is fluid. A self-correction backend can be viewed
as a remasking sampler with a particular trained-in correction head; a
remasking sampler with a strong supervision signal on confidence (RemeDi,
CDLM) can be viewed as a partial self-correction backend. The thesis
treats them as distinct because the timing question is different: a
remasking sampler's correction frequency is a sampler hyperparameter
integrated into the entire trajectory, whereas a self-correction backend
supports a per-step decision to invoke the refinement head.

\section{Uncertainty-driven decoding}
\label{lit:sec:uncertainty}

A second family of inference-time mechanisms uses model uncertainty to
drive decoding decisions. \textcite{benhamu2025ebsampler} introduce
entropy-bounded unmasking: at each step, the sampler unmasks the largest
set of positions whose total predicted entropy stays below a target. The
rationale is that unmasking many low-entropy positions in parallel does not
noticeably degrade quality, while unmasking high-entropy positions in
parallel does. The method accelerates inference at fixed quality and
exposes a target-entropy hyperparameter that can itself be tuned per step.

\textcite{chen2025denoisingentropy} formalise the role of entropy as a
per-step quality signal, introduce a denoising-entropy quantity that
decomposes total entropy into reducible and irreducible parts, and show
that decoding paths that prioritise reducing entropy at the right rate
produce better samples. The method optimises an entropy-trajectory
objective and is most directly related to the corrector-timing question
of this thesis: it is a per-step scheduling problem, although its decision
variable is the unmasking policy rather than corrector placement.

These uncertainty-driven mechanisms supply candidate signals for corrector
timing: per-step entropy, per-step inverse margin, and entropy-trajectory
derivatives are all observable quantities that could serve as proxies for
the marginal value of correction at a given step. They motivate the
marginal-ranker baseline developed in Chapter~\ref{ch:contribution}.

\section{Informed correctors for discrete diffusion}
\label{lit:sec:informed-correctors}

\textcite{zhao2024informedcorrectors} introduce informed correctors for
discrete diffusion models. An informed corrector is a Markov transition
designed to refine the current state of a discrete diffusion sampler using
model-derived information beyond the trained denoiser's per-position
distribution. The two primary instances developed in their paper are based
on locally balanced proposals \parencite{zanella2020informedproposals} and
on confidence-weighted resampling.

The first family of correctors in \textcite{zhao2024informedcorrectors}
uses locally balanced proposals in the sense of
\textcite{zanella2020informedproposals}: at each corrector step a small
number of positions are selected for resampling, and a proposal
constructed from a balancing function applied to local probability ratios
is used to update those positions. The second family uses
confidence-based resampling, in which positions with low predicted
confidence are selected for resampling. The precise proposal form,
target, and acceptance rule for each variant are stated in the notation
of the original paper. Both correctors require a notion of action
set---the positions on which the corrector is allowed to operate---and a
notion of corrector budget---how many corrector steps to take.

The relevance for the present thesis is that informed correctors
instantiate a structural pattern: a fixed Markov transition with a
designated action set, a per-application cost, and an empirical effect on
sample quality. The thesis takes this pattern as given and asks the next
question: with the corrector fixed and a budget for placements, where
should the corrector be applied? This question is partly answered by the
literature on adaptive samplers and partly left open. The remainder of this
chapter surveys the most relevant adjacent work and then turns to the
timing question.

\section{MCMC background: informed proposals, Gibbs samplers, and
convergence}
\label{lit:sec:mcmc}

Informed correctors borrow heavily from the MCMC literature on discrete
state spaces, and several theoretical strands are worth situating because
they recur in the thesis.

\textcite{zanella2020informedproposals} introduces locally informed
proposals for discrete-space MCMC, parameterising the proposal
distribution via a balancing function $g:[0,\infty)\to[0,\infty)$
applied to local probability ratios. The framework includes the Barker
proposal and the locally balanced proposals as members of the family,
and proves asymptotic optimality of the locally balanced choice under
specified high-dimensional regimes. The proposals can be combined with a
Metropolis--Hastings accept/reject step or used as drop-in samplers
depending on the target.

The Gibbs sampler is the canonical block-Markov-chain Monte Carlo
procedure: at each step, one selects a coordinate (or block of coordinates)
and updates it from its full conditional distribution. Theoretical results
on Gibbs samplers go back to \textcite{roberts1997updating}, who study how
the choice of update scheme (random scan, systematic scan, blocked Gibbs)
affects mixing under specific covariance structures. The discrete Gibbs
sampler is also the canonical model for masked-diffusion correctors: a
corrector that resamples the marginal at a single position from the
predicted distribution conditional on the rest is precisely a Gibbs step,
and a corrector that resamples a small block is a blocked Gibbs step.

\textcite{ascolani2024entropycontraction} establish entropy-contraction
results for Gibbs samplers under log-concavity assumptions.
\textcite{secchi2025mwggap} analyse the spectral gap of
Metropolis-within-Gibbs samplers under log-concavity. Both works give
quantitative convergence rates that depend on specific structural
properties of the target (curvature, conditioning). For discrete
diffusion correctors the targets are typically not log-concave in any
natural geometry, and these results do not transfer directly. They
nonetheless supply a vocabulary for thinking about why some corrector
steps are more valuable than others: corrector placements at points
where the chain is far from the target make more progress per step than
placements at points where the chain is already close.

The classical comparison theorems for reversible Markov chains
\parencite{diaconis1993comparison} provide one further piece of vocabulary.
They allow one to bound the spectral gap of one Markov kernel in terms of
another, with a constant depending on a Radon--Nikodym-type ratio. In the
corrector context, this is the right framework for asking whether two
correctors with different action sets or different acceptance rules have
comparable mixing properties. The thesis does not develop this line of
analysis directly, but it draws on the spectral-gap vocabulary in the
diagnostic framework of Chapter~\ref{ch:contribution} when comparing
pairwise-aware to marginal schedules.

\section{ProSeCo: self-correcting masked diffusion}
\label{lit:sec:proseco}

\textcite{schiff2026proseco} introduce ProSeCo, a self-correcting masked
diffusion model that jointly trains an unmasking head and a correction
head. The model has the same overall architecture as a standard MDLM but
is trained on a curriculum that includes both standard mask-prediction
examples and synthetic `mistake' examples in which the supervision target
differs from the model's previously committed token. The result is a
single network that supports two operations at inference: ordinary
unmasking, used to fill in masked positions, and refinement, used to
revise committed positions.

The refinement head defines an informed corrector with a specific action
set: at any step $t$, the set of revisable positions is the set of
currently committed positions, and the corrector resamples those positions
from the refinement head's predicted distribution. The corrector cost is
one additional forward pass per refinement step.
\textcite{schiff2026proseco} report that ProSeCo improves sample quality
on standard benchmarks and exhibits the qualitative behaviour expected of
a self-correction model: when refinement is enabled, committed tokens that
conflict with later-revealed context tend to be revised toward the data
distribution.

The relevance for the thesis is twofold. First, ProSeCo supplies a concrete
informed corrector with a well-defined action set, a measurable cost per
call, and an empirical track record of helpful behaviour. Second,
ProSeCo's predictor--corrector structure is the closest masked-diffusion
analogue of the continuous predictor--corrector samplers of
\S\ref{lit:sec:pred-corr}: the predictor advances the trajectory, the
corrector refines without changing the time index, and the inference-time
cost is partitioned between the two. The thesis takes ProSeCo as the
empirical backend and studies the corrector-timing question on top of its
trained refinement head, holding the head and its training fixed.

ProSeCo is one instance of a broader family of self-correcting and
corrective masked-diffusion designs. \textcite{huang2025remedi} jointly
train the denoiser with a confidence head used by the sampler to decide
what to remask. \textcite{zhang2026cdlm} explicitly train for corrective
behaviour by augmenting the loss with examples that require the model to
assign lower confidence to incorrect tokens. \textcite{wu2026dsl} train a
single denoiser to be SNR-invariant across mask schedules, supporting
remasking-induced trajectories without retraining. These approaches differ
in whether they expose a separate refinement head (ProSeCo) or fold the
corrective behaviour into the unmasking head (RemeDi, CDLM, DSL), and in
whether the corrector is invoked through an explicit per-step decision
(ProSeCo) or as part of the standard sampler (RemeDi).

\section{PRISM: quality-guided refinement}
\label{lit:sec:prism}

\textcite{kim2025prism} introduce PRISM, a method for self-correction that
does not require retraining the base denoiser. PRISM trains a separate
per-token quality scorer on synthetic mistake examples and uses the
resulting scores at inference to identify low-quality committed tokens,
remask them, and let the base denoiser regenerate them. The method is
plug-in in the sense that any pretrained masked diffusion language model
can be paired with a separately trained PRISM scorer.

PRISM has two distinct uses that should be distinguished, because the
literature occasionally conflates them and the present thesis takes a
careful position on each.

First, \emph{PRISM as a non-separable quality-guided refinement
mechanism}. In its native deployment, PRISM uses the full vector of
per-token quality scores to choose a refinement set, so the refinement
is not equivalent to ranking time steps by a single scalar. Used this
way, PRISM is an informed-refinement backend in its own right. The
present thesis does not test this deployment of PRISM and does not rule
it out as a backend.

Second, \emph{PRISM-style quality mass as a per-step separable timing
signal}. The aggregate quality mass at step $t$, which the thesis denotes
$QM_t$, is a single scalar that summarises how confident the PRISM scorer
is in the currently committed tokens. As a per-step signal, $QM_t$ is in
the same class as $H_t$ (entropy) and $M_t^{-1}$ (inverse margin)---a
separable score on which marginal ranking can be evaluated. The
contribution chapters of this thesis test marginal ranking on signals of
this class. A negative result for tested separable rankers does not rule
out PRISM as a non-separable backend; it limits PRISM-style quality mass
when used as a separable timing score in the tested ranker class.

This distinction matters because the literature on quality-guided
refinement is closely connected to the literature on corrector timing
without being identical. Quality scores are useful as signals for
refinement; the question of where in the trajectory to spend a fixed
number of refinement calls remains open even when the refinement mechanism
itself has been validated.

\section{Trajectory search and test-time scaling}
\label{lit:sec:search}

A separate line of work spends additional inference compute on searching
over trajectories or candidates. \textcite{lee2025iterref} introduce
IterRef, an iterative-refinement procedure for masked text-to-mask models
that uses reward signals to guide each refinement round.
\textcite{bilal2026s3} introduce S$^3$, a stratified scaling search that
branches over candidate paths at each step and keeps the best according
to a verifier. \textcite{dang2025particlegibbs} introduce PG-DLM, a
particle Gibbs sampler for diffusion language models that combines
sequential Monte Carlo with MCMC to refine entire trajectories using
reward feedback. \textcite{koh2024eags} introduce entropy-adaptive Gibbs
sampling, in which Gibbs steps are scheduled in proportion to estimated
entropy and refinement is restricted to high-entropy regions of the
sequence. \textcite{deschenaux2026psisamplers} develop a curriculum-based
corrector family ($\Psi$-samplers) that combines several elements of
corrector kernel design with schedule curriculum learning.
\textcite{hasan2026fkc} introduce discrete Feynman--Kac correctors, an
SMC-style approximation of the reverse process in which a population of
weighted candidate trajectories is propagated and reweighted at each
step.

These approaches differ from the thesis question in two structural
respects. First, they search over candidate trajectories or outputs; they
do not, in general, fix a single trajectory and ask only at which time
steps an additional corrector should be applied. Second, they use external
feedback---reward, verifier, ancestor signal, or explicit schedule-gain
estimation---that is not generally available in a deployable inference-time
scheduler. The thesis treats search procedures with true-feedback access
as structural existence results: they can show that good schedules exist
(because, given true feedback, they can find them), but they do not
establish that an inference-time scheduler can be learned from observable
signals alone.

\section{From informed correctors to timing}
\label{lit:sec:to-timing}

The methods reviewed in this chapter all modify inference-time behaviour,
but they do so at different layers of the masked-diffusion pipeline. The
base model supplies a denoiser. Remasking samplers change the trajectory.
Self-correction backends supply a fixed refinement transition with a
designated action set. Quality-guided refinement biases the inference
distribution. Trajectory search branches over candidates. The thesis
question---corrector timing---takes one specific configuration: a base
masked diffusion model, a fixed informed corrector, a fixed quality
functional, and a fixed corrector-placement budget. The remaining degree
of freedom is the schedule: the set of denoising time steps at which the
corrector is applied.

This timing question has analogues in several adjacent literatures, and it
is worth previewing the connections. Set-function optimisation and
submodularity supply the language for offline schedule selection: with a
corrector budget $B$, choosing a size-$B$ schedule that maximises a
quality functional is a cardinality-constrained set-function problem, and
the theory of submodular optimisation classifies when greedy selection is
provably near-optimal. The classical result is due to
\textcite{nemhauser1978submodular}---greedy maximisation of a monotone
submodular function under a cardinality constraint achieves a $(1-1/e)$
approximation---and adaptive submodularity extends the argument to
policies that observe outcomes between actions
\parencite{golovin2011adaptive}. These results are not directly applicable
here because corrector schedule gains are not, in general, monotone
submodular, but they supply a useful classification scheme: additive,
submodular, complementary, or higher-order.

Sequential Monte Carlo and particle filtering supply the language for
trajectory-level refinement. The SMC pattern---propagate a population of
trajectories, monitor degeneracy, and rejuvenate when needed---maps
directly onto the masked-diffusion corrector: propagate a partially masked
trajectory, monitor a degeneracy diagnostic (entropy, margin, quality
mass), and apply a corrector when the diagnostic exceeds a threshold.
This analogy motivates uncertainty-triggered online policies and is the
closest non-trivial parallel for the corrector-timing question.
\textcite{hasan2026fkc} and \textcite{dang2025particlegibbs} make the
particle/SMC connection explicit.

Sensor scheduling and active sensing supply an analogous language for
resource-constrained measurement allocation in linear-Gaussian state
estimation. The discrete masked-diffusion setting does not satisfy the
linear-Gaussian assumption, so these criteria transfer as design
intuition rather than as importable bounds, but the underlying question
is structurally similar: when should a scarce inference-time call be
spent to reduce a downstream error?

Value-of-information arguments suggest that a corrector call should be
spent when its expected benefit exceeds the opportunity cost of saving
the budget for a future step. Index policies for resource-constrained
allocation give cheap policy classes for this kind of budget-aware
online decision, with the caveat that their optimality typically rests
on separability assumptions that may or may not hold in the
masked-diffusion setting.

These analogies are intellectually fertile but not, on their own,
decisive. None of them solves the corrector-timing problem off the shelf.
They do, however, suggest the structure that a thesis-grade answer should
have. With a fixed predictor, a fixed corrector, and a fixed budget, one
can ask at three levels of generality. At the coarsest level, is there any
schedule that beats uniform placement? At the next level, can a separable
per-step ranker recover most of the available headroom? At the next, do
schedule interactions matter, and if so are they captured by pairwise
terms or do they require higher-order or search-based scheduling? The
contribution chapters of this thesis formalise these questions through
the schedule-gain object $G(S)$, the additive surrogate $A(S)$, and the
pairwise surrogate $Q(S)$, and develop diagnostics for classifying a
given backend into the corresponding regimes (no-op, marginal,
interaction-driven, higher-order/search-driven, or online-decision). The
empirical chapters then ask which of these regimes a specific
masked-diffusion backend occupies.

The chapter has surveyed the literature that supplies the corrector
kernel, the action set, the candidate signals, and the structural
vocabulary. The remaining task is to fix these ingredients and to develop
a principled framework for corrector timing. That is the subject of
Chapter~\ref{ch:contribution}.

%% =====================================================================
%% End of literature-review draft (chapters 2--5)
%% =====================================================================
